\section*{Abstract}

The manual evaluation of Ground Penetrating Radar (GPR) data is reaching its limits due to the constantly growing amount of data. An automated solution for segmenting the radargrams is therefore required to ensure efficient analysis. In this work, unsupervised machine learning methods from seismics and seismology are transferred to GPR data to enable automated segmentation. A Self-Organizing Map (SOM) is used as the central tool. First, radargrams from an archaeological site at Lüchow were preprocessed, and eleven different features were extracted from them. Subsequently, Self-Organizing Maps were trained using different feature combinations of a radargram in order to achieve meaningful clustering. The transferability of the model was tested by applying it to an independent dataset. \\
The studies show that four of the features—instant amplitude, quadrature, instantaneous frequency, and kurtosis—are particularly well suited for segmentation, and that SOM is capable of forming meaningful clusters from the calculated features, thus distinguishing heterogeneous from homogeneous areas and revealing noise. Furthermore, it was successfully demonstrated that the model trained on one radargram also delivers results that are applicable to other radargrams, confirming the robustness and transferability of the approach.\\
In summary, self-organizing maps represent an effective method for automatically segmenting radargrams and can make a valuable contribution to increasing the efficiency of GPR data ana\-lysis.

\clearpage

\section*{Zusammenfassung}

Die manuelle Auswertung von Ground Penetrating Radar (GPR) Daten, stößt aufgrund stetig wachsender Datenmengen an ihre Grenzen. Eine automatisierte Lösung zur Segmentierung von Radargrammen ist daher erforderlich, um eine effiziente Analyse zu gewährleisten. In dieser Arbeit werden Methoden des unüberwachten maschinellen Lernens aus Seismik und Seismologie auf GPR-Daten übertragen, um eine automatisierte Segmentierung zu ermöglichen. Dabei wird eine Self-Organizing Map (SOM) als zentrales Werkzeug eingesetzt. Zunächst wurden Radargramme einer archäologischen Stätte bei Lüchow vorverarbeitet und aus diesen elf verschiedene Merkmale extrahiert. Anschließend wurden Self-Organizing Maps mithilfe unterschiedlicher Merkmalskombinationen eines Radargramms trainiert, um eine aussagekräftige Clusterbildung zu erreichen. Die Übertragbarkeit des Modells wurde durch die Anwendung auf einen unabhängigen Datensatz geprüft. \\
Die Untersuchungen zeigen, dass die vier Merkmale Instantaneous Amplitude, Quadrature, Instantaneous Frequency und Kurtosis besonders gut für die Segmentierung geeignet sind und dass die SOM in der Lage ist, aus den berechneten Merkmalen sinnvolle Cluster zu bilden, um damit heterogene von homogenen Bereichen zu unterscheiden und Rauschen aufzuzeigen. Darüber hinaus konnte erfolgreich demonstriert werden, dass das auf einem Radargramm trainierte Modell auch auf anderen Radargrammen anwendbare Ergebnisse liefert, was die Robustheit und Übertragbarkeit des Ansatzes bestätigt.\\
Zusammenfassend lässt sich festhalten, dass Self-Organizing Maps eine effektive Methode zur automatischen Segmentierung von Radargrammen darstellen und einen wertvollen Beitrag zur Effizienzsteigerung bei der GPR-Datenauswertung leisten können.